\chapter*{Introducción}
\addcontentsline{toc}{chapter}{Introducción}
\markboth{Introducción}{Introducción}

La bioenergía es la fuente renovable con mayor aporte al suministro energético mundial. De acuerdo con la Asociación Mundial de la Bioenergía (WBA), en 2023 aportó aproximadamente 56~EJ de un total mundial de 622~EJ (alrededor del 9~\%); dado que las energías renovables en conjunto sumaron 92~EJ, cerca de seis de cada diez unidades de energía renovable provinieron de la biomasa \cite{wba2025gbsr}. Este predominio obedece principalmente al uso térmico (cocción, calefacción y calor de proceso industrial), diferenciándose entre el aprovechamiento tradicional de baja eficiencia y la bioenergía moderna en instalaciones industriales controladas.

En la generación eléctrica, su participación es más acotada, pero presenta un rasgo distintivo frente a la energía solar fotovoltaica y la eólica: es una fuente \textbf{despachable}, ya que el combustible puede almacenarse y la central puede generar según las necesidades del sistema eléctrico. Al cierre de 2025 la capacidad mundial instalada de bioenergía alcanzó 154~GW (cerca del 3~\% de los 5\,149~GW de capacidad renovable total), con una adición anual de 3,4~GW frente a 692~GW del conjunto de renovables \cite{irena2026capacity}. En 2024 produjo 698~TWh (7,1~\% de la electricidad renovable) con 151~GW instalados \cite{wba2025gbsr}, lo que arroja un factor de planta medio cercano al 53~\%:
\begin{equation*}
\mathrm{FP} = \frac{698\text{ TWh}}{151\text{ GW} \times 8\,760\text{ h}} = \frac{698\,000\text{ GWh}}{1\,322\,760\text{ GWh}} \approx 0{,}53
\end{equation*}

\begin{table}[h!]
\centering
\footnotesize
\caption{Indicadores globales de la bioenergía (2023--2025).}
\label{tab:biomasa_cifras}
\begin{tabularx}{\textwidth}{@{}X l l@{}}
\toprule
\textbf{Indicador} & \textbf{Valor} & \textbf{Fuente} \\
\midrule
Suministro mundial de bioenergía en Panama (2023) & 56~EJ ($\sim$9~\% total; $\sim$60~\% renovables) & \cite{wba2025gbsr} \\
Electricidad generada con bioenergía (2024) & 698~TWh (7,1~\% de la generación renovable) & \cite{wba2025gbsr} \\
Capacidad instalada de bioenergía (fin 2025) & 154~GW (+3,4~GW en 2025; +2,3~\%) & \cite{irena2026capacity} \\
Capacidad renovable total (fin 2025) & 5\,149~GW (+692~GW en 2025; +15,5~\%) & \cite{irena2026capacity} \\
Principales productores eléctricos (2024) & China: 209~TWh; Brasil: 58~TWh; EE.\,UU.: 47~TWh & \cite{wba2025gbsr} \\
Variación de costos en 2024 frente a 2023 & Costo instalado +16~\%; LCOE +13~\% & \cite{irena2025costs} \\
\bottomrule
\end{tabularx}
\end{table}

La presente monografía aborda la energía de biomasa desde la perspectiva de las plantas de potencia. El Capítulo~1 desarrolla los fundamentos fisicoquímicos, clasificación, propiedades energéticas y logística del combustible. El Capítulo~2 analiza las rutas y tecnologías de conversión: combustión directa, gasificación, pirólisis y digestión anaeróbica. El Capítulo~3 examina las escalas y configuraciones de centrales en el mundo, casos de referencia y su aporte a la estabilidad del sistema eléctrico. Finalmente, el Capítulo~4 evalúa la economía de generación, las emisiones de ciclo de vida, la captura de carbono (BECCS) y el panorama en Panamá, seguido de las conclusiones generales.

\clearpage
